\subsection{By the end of this lecture}
\begin{itemize}
    \item You should have an overview of super-twisting sliding mode control, strengths, weaknesses, variations and solutions
\end{itemize}

\subsection{Sliding mode recap}
Second-order system 
\begin{equation*}
    \begin{split}
        \dot{x}_1 &= x_2 \\
        \dot{x}_2 &= h(x) + g(x)u
    \end{split}
\end{equation*}
where $h$ and $g$ are unknown nonlinear functions of time and $g(x) \geq g_0 > 0, \forall x$. If we want to stabilize the origin we can choose 
\begin{equation*}
    s = \lambda x_1 + x_2
\end{equation*}
with derivative 
\begin{equation*}
    \dot{s} = \lambda x_2 + h(x) + g(x)u
\end{equation*}
if $s = 0$ we obtain our sliding manifold (surface). This gives the dynamics
\begin{equation*}
    x_2 = -\lambda x_1
\end{equation*}
Meaning that the dynamics go exponentially to zero on the sliding manifold.

Therefore, we want to choose a control law that sends the trajectory to the sliding manifold and stays there, i.e. $s = 0$ and $\dot{s} = 0$. Through Lyapunov analysis we end up with the function
\begin{equation*}
    u = -\beta(x) sgn (s)
\end{equation*}
where 
\begin{equation*}
    sgn(s) = 
    \begin{cases}
        1, &s > 0 \\
        0, &s = 0 \\
        -1, &s < 0 \\
    \end{cases}
\end{equation*}